% ====================================================================
\section{The computational model: monotone Boolean straight-line programs}\label{sec:model}
% ====================================================================

\begin{definition}[Monotone Boolean straight-line program]\label{def:mslp}
A \emph{monotone Boolean straight-line program} (MSLP) is a tuple
$\mathcal P=(r,\iota_0,(\pi_1,\dots,\pi_L),o)$ where:
\begin{itemize}
\item $r\ge 1$ is the number of \emph{registers} $R_1,\dots,R_r$, each holding one bit;
\item $\iota_0\in\{0,1\}^r$ is the \emph{initial content} of the registers;
\item each \emph{instruction} $\pi_\ell$ has the form
  $R_{d_\ell}\leftarrow R_{a_\ell}\circ_\ell R_{b_\ell}$, with $d_\ell,a_\ell,b_\ell\in[1,r]$ and
  $\circ_\ell\in\{\wedge,\vee\}$;
\item $o\in[1,r]$ is the \emph{designated register}.
\end{itemize}
The program is encoded by listing $r$, $\iota_0$, and the quadruples $(d_\ell,\circ_\ell,a_\ell,b_\ell)$
with indices in binary.
\end{definition}

\begin{definition}[Semantics]\label{def:semantics}
The \emph{register states} $\rho_0,\dots,\rho_L:[1,r]\to\{0,1\}$ are $\rho_0=\iota_0$ and
\[
\rho_\ell(x)=\begin{cases}\rho_{\ell-1}(a_\ell)\circ_\ell\rho_{\ell-1}(b_\ell)&\text{if }x=d_\ell,\\
\rho_{\ell-1}(x)&\text{otherwise.}\end{cases}
\]
The \emph{value} of $\mathcal P$ is $\val(\mathcal P)=\rho_L(o)$.
\end{definition}

The model has the four features required by the reduction:
\begin{itemize}
\item \emph{Boolean registers.} The only data are bits.
\item \emph{Monotone instructions.} Only $\wedge$ and $\vee$ occur. Constants enter only through
  $\iota_0$, and there is no negation.
\item \emph{Reuse of values.} A register may be read by any number of later instructions, and may
  also be overwritten.
\item \emph{No conditional jumps.} The instructions are executed once each, in their listed order.
  Which registers are read and written at step $\ell$ is fixed by the program text, independently of
  the data.
\end{itemize}
We call the last property \emph{obliviousness}. It is what allows a fixed physical schedule to
execute the program, as in \cref{sec:processor}.

\begin{definition}[Single-assignment normal form]\label{def:normalform}
An MSLP in \emph{normal form} has registers $R_1,\dots,R_m$ with $m\ge 2$. There is a set
$I\subseteq[1,m-1]$ of \emph{input registers}, and instructions indexed by the
\emph{instruction registers} $k\in[1,m]\setminus I$, executed in increasing order of $k$. Each
instruction is $R_k\leftarrow R_{i_k}\circ_k R_{j_k}$ with $i_k,j_k<k$; the case $i_k=j_k$ is
allowed. Thus every register is written at most once, and only after all the registers it reads.
The designated register is $R_m$, and the last instruction is $R_m\leftarrow R_o\wedge R_o$ for
some $o<m$. Input registers have initial values $\iota_0(i)$. Instruction registers have initial
value $0$, which is never read. Register $R_1$ is necessarily an input, because an instruction $k$
needs operands smaller than $k$. In particular every instruction register satisfies $k\ge 2$.
\end{definition}

A program in normal form is the same thing as a topologically ordered monotone circuit: constants
are the input registers, gates are the instruction registers, and the predecessors of gate $k$ are
$i_k,j_k<k$. Arbitrary fan-out corresponds to repeated reads.

\begin{lemma}[Normal form in logspace]\label{lem:normalform}
There is a logspace transducer that maps every MSLP $\mathcal P$ to an MSLP $\mathcal P'$ in normal
form with $\val(\mathcal P')=\val(\mathcal P)$ and $m=r+L+1$ registers.
\end{lemma}

\begin{proof}
Let $\mathcal P=(r,\iota_0,(\pi_\ell)_{\ell\le L},o)$. For $x\in[1,r]$ and $1\le\ell\le L+1$ define
\[
\lambda(x,\ell)=\begin{cases}r+\max\{\ell'<\ell:\ d_{\ell'}=x\}&\text{if this set is nonempty},\\
x&\text{otherwise}.\end{cases}
\]
Thus $\lambda(x,\ell)$ names the register of $\mathcal P'$ that holds the content of $R_x$ just
before instruction $\ell$ of $\mathcal P$. Let $\mathcal P'$ have registers $1,\dots,r+L+1$, inputs
$I=[1,r]$ with the values $\iota_0$, and for $1\le\ell\le L$ the instruction
\[
R'_{r+\ell}\leftarrow R'_{\lambda(a_\ell,\ell)}\circ_\ell R'_{\lambda(b_\ell,\ell)},
\]
followed by $R'_{r+L+1}\leftarrow R'_{\lambda(o,L+1)}\wedge R'_{\lambda(o,L+1)}$. All operand indices
are smaller than the index written, so $\mathcal P'$ is in normal form.

By induction on $\ell$, the value of $R'_{r+\ell}$ equals $\rho_\ell(d_\ell)$, and
$\rho_{\ell}(x)$ equals the value of $R'_{\lambda(x,\ell+1)}$ for every $x$. The base case is
$\lambda(x,1)=x$. For the inductive step, if $x\ne d_\ell$ then
$\lambda(x,\ell+1)=\lambda(x,\ell)$, and if $x=d_\ell$ then $\lambda(x,\ell+1)=r+\ell$. Hence
$\val(\mathcal P')=\rho_L(o)\wedge\rho_L(o)=\val(\mathcal P)$.

For logspace, each output symbol is computed by a loop over $\ell'$ from $\ell-1$ down to $1$
that compares $d_{\ell'}$ with $x$ by rereading the input. Only a constant number of indices of
$O(\log(r+L))$ bits are stored.
\end{proof}

\begin{remark}\label{rem:finaland}
The final instruction $R_m\leftarrow R_o\wedge R_o$ copies the designated value. It makes the last
instruction an $\AND$, which the eventual-activation construction of \cref{sec:ever} uses.
For a topologically ordered circuit with designated output $o$, appending this gate is the
convention of the original manuscript.
\end{remark}
